% !PRINTED 124
\weberpage{124}
%<<CONT

\gesperrt{wird, dass also die untere Grenze auch hier ein Minimum ist.}

Die \gesperrt{untere Grenze} von $X$ in irgend einem Theile des Bereiches ist entweder gleich $g$ oder grösser als $g$.

Ein Grössengebiet, das durch die Ungleichheitsbedingungen
\begin{equation}
 |x|\le R,\qquad -R\le y\le \alpha
\tag{7}
\end{equation}
bestimmt ist, wird in unserer Fig. 2 durch das Segment $(PQ\alpha Q')$, das wir das Segment $(P,\alpha)$ nennen wollen, dargestellt.

% Exact 300-dpi crop from the registered Volume-I scan, printed Fig. 2.
\weberfigureimage{weber-b1-fig02.png}{0.62}

Wir bestimmen nun zunächst einen Werth $\eta$ von $y$ durch einen Schnitt $(\mathfrak A,\mathfrak B)$ folgendermassen: Eine Zahl $\alpha$ zwischen $-R$ und $+R$ wird in $\mathfrak A$ aufgenommen, wenn die \gesperrt{untere Grenze} von $X$ in dem Segment $(P,\alpha)$ grösser als $g$ ist, und ein Werth $\beta$ wird in $\mathfrak B$ aufgenommen, wenn die \gesperrt{untere Grenze} von $X$ in dem Segment $(P,\beta)$ gleich $g$ ist. Dieser Schnitt $(\mathfrak A,\mathfrak B)$ definirt eine Zahl $\eta$ von der Eigenschaft, dass die \gesperrt{untere Grenze} von $X$ in dem Bereich $(P,y)$, wenn $y<\eta$ ist, grösser als $g$, und wenn $y>\eta$ ist, gleich $g$ ist.

Nun ist nach \S.\,31, VI
\[
 |f(\eta+iz)|
\]
eine stetige Function von $z$ in dem Intervall
\begin{equation}
 -\sqrt{R^2-\eta^2}\le z\le \sqrt{R^2-\eta^2},
\tag{8}
\end{equation}
was in der Fig. 2 durch $M,M'$ bezeichnet ist.

Diese Function erhält also nach 3. für irgend einen Werth $\zeta$ von $z$ in diesem Intervall einen Minimumwerth $\gamma$, so dass, wenn
\[
 \xi=\eta+i\zeta
\]
gesetzt wird,
\begin{equation}
 |f(\xi)|=\gamma
\tag{9}
\end{equation}
wird.
